Nach dem Stefan-Boltzmann-Gesetz (FoTaBe S.~172) strahlt eine Oberfläche $A$ mit der Emissionszahl $\varepsilon$ und der absoluten Temperatur $T$ die Leistung $P = \varepsilon\,\sigma\,A\,T^4$ ab. Für poliertes Aluminium ist $\varepsilon = \eAlO$ (FoTaBe S.~190), $\sigma = \sigO$ (FoTaBe, Umschlag hinten), und $T = \TK$:
\begin{align}
 P &= \PKF \\
   &= \eAlO\cdot\sig\cdot4\pi\cdot\left(\rK\right)^2\cdot\left(\TK\right)^4 \\
   &= \PK \approx \result{\PKP{0}{2}}
\end{align}
Die Kugel nimmt gleichzeitig Strahlung aus der Umgebung auf; netto verliert sie deshalb etwas weniger Leistung. Wegen der kleinen Emissionszahl strahlt poliertes Aluminium nur wenig ab, eine schwarze Kugel gäbe bei gleicher Temperatur ein Vielfaches ab.
